\chapter{PROCJENA SVOJSTAVA OKOLINE}
\label{ch:procjena}

Zadatak traži da robot svojstva okoline procjenjuje iz mjerenja sile i momenta tijekom
interakcije, a ne da mu se ona zadaju unaprijed. Ovo poglavlje opisuje kako se iz veličina
dostupnih na robotu dolazi do sile i momenta na vrhu alata, te kako se iz njihova odnosa prema
gibanju u načelu procjenjuju svojstva okoline poput krutosti i ograničenja gibanja.

Naglasak je na tome da su sve procjene izvedene isključivo iz veličina koje bi imao i stvarni
robot (poze vrha alata i procijenjene sile) a nikada iz stanja zglobova vrata, koje
simulator doduše zna, ali stvarni robot ne bi. Stanje zgloba vrata koristi se samo kao
referentna vrijednost pri provjeri procjenitelja, nikada u upravljačkoj petlji.

\section{Procjena sile i momenta na vrhu alata}
\label{sec:wrench}

Manipulator nema zaseban senzor sile na zapešću, nego se sila i moment na vrhu alata procjenjuju
iz momenata u zglobovima. Veza je posljedica načela virtualnog rada: za manipulator u ravnoteži,
moment u zglobovima koji uravnotežuje vanjsku silu i moment na vrhu alata dan je izrazom

\begin{equation}
    \bm{\tau} = \bm{J}^{\mathsf{T}}(\bm{q})\,\bm{w} ,
    \label{eq:jt}
\end{equation}

gdje je $\bm{\tau}$ vektor momenata u zglobovima, $\bm{q}$ vektor zglobnih koordinata,
$\bm{J}(\bm{q})$ Jakobijan manipulatora, a $\bm{w} = [\,\bm{F}^{\mathsf{T}} \ \
        \bm{M}^{\mathsf{T}}\,]^{\mathsf{T}}$ vektor sile i momenta na vrhu alata (engl. \textit{wrench}).
Procjena sile i momenta dobiva se rješavanjem tog sustava po $\bm{w}$.

Za manipulator sa sedam zglobova i šesterodimenzijskim vektorom sile i momenta sustav
\eqref{eq:jt} je preodređen, pa se rješava u smislu najmanjih kvadrata pomoću
Moore-Penroseova pseudoinverza transponiranog Jakobijana:

\begin{equation}
    \hat{\bm{w}} = \left( \bm{J}^{\mathsf{T}} \right)^{+} \bm{\tau} .
    \label{eq:wrench-est}
\end{equation}

Jakobijan se računa iz trenutne konfiguracije manipulatora, dobivene iz stabla koordinatnih
sustava, pa procjena ne zahtijeva nikakav dodatan senzor osim već postojećih mjerila momenta u
zglobovima.

\subsection{Svojstva i ograničenja procjene}
\label{sec:wrench-svojstva}

Izraz \eqref{eq:wrench-est} pretpostavlja da je izmjereni moment u zglobovima posljedica
isključivo sile i momenta na vrhu alata. U stvarnosti mjereni moment sadrži i doprinos
vlastite težine manipulatora i prihvatnice, koji o kontaktu ništa ne govori.
\texttt{tcp\_wrench\_estimator} tu težinu ne kompenzira proračunom, nego je uklanja umjerivanjem:
neposredno prije zanimljivog dijela zadatka, dok manipulator miruje u konfiguraciji bliskoj
onoj u kojoj će se odvijati kontakt i dok kontakta još nema, očitani se moment pohrani kao
poravnanje na nulu (\textit{tare}) i oduzima od svih narednih očitanja. Umjeravanje je stoga
vezano uz konfiguraciju: ono uklanja gravitacijski doprinos onoj pozi u kojoj je izvedeno, a
odstupanje raste kako manipulator odmiče od te poze u sljedećim koracima gibanja.

Ni nakon umjeravanja procjena nije bez pogreške. U slobodnom gibanju, bez kontakta, procjena
sile pokazuje šum reda pola njutna, što je razina ispod koje se kontakt ne može pouzdano
razlučiti od mjerne nesigurnosti. Procjena momenta uz to zadržava sustavni rezidual reda
desetak njutnmetara, koji potječe od preostale, nepotpuno uklonjene dinamike manipulatora,
ponajviše iz odstupanja stvarne poze od one u kojoj je umjeravanje provedeno. Na tipičnom kraku
hvata taj rezidual odgovara sili reda dvadesetak njutna.

Posljedica je da procjena dobro razlučuje smjer i prisutnost kontaktne sile, ali njezin
apsolutni iznos treba uzimati s obzirom na navedeni rezidual, a umjeravanje treba ponoviti
kad god se manipulator znatnije udalji od poze u kojoj je posljednje umjeravanje provedeno.
Ponovno umjeravanje tijekom zadatka provodi se uvjetno, samo kad je trenutno očitana sila već
niska. U suprotnom bi se u nulu upisala stvarna kontaktna sila, umjesto samo preostalog
gravitacijskog doprinosa. Za upravljanje u ovom radu to je dovoljno, jer se sila koristi za
otkrivanje kontakta i za ograničavanje njezina iznosa, a ne kao precizno mjerilo.

\section{Ispitivanje ograničenja gibanja}
\label{sec:constraint}

Iz procijenjene sile i istovremeno mjerenog gibanja može se utvrditi u kojim je smjerovima gibanje
vrata slobodno, a u kojima ograničeno. Postupak je izveden kao zasebna rutina
(\texttt{door\_probe.py}) koja se pokreće iz uspostavljenog hvata. Vrh alata izvodi male,
silom nadzirane pomake u šest smjerova, uzduž triju osi koordinatnog sustava prihvatnice u oba
smjera, po četiri koraka od 3 mm.

\subsection{Mjera koja razlikuje slobodan smjer od ograničenog}

Kriterij nije iznos sile, već koeficijent sile po naređenom pomaku, izražen u njutnima po metru.
Razlika između dvaju slučajeva jasno je vidljiva u obliku niza očitanja. U slobodnom smjeru sila
naraste do razine potrebne da se svlada trenje u zglobu vrata i ondje se zadrži, jer se vrata
gibaju. U ograničenom smjeru ništa se ne giba, pa se lanac između platforme i vrata samo napinje
i sila monotono raste. Smjer se proglašava slobodnim ispod 8 N/mm, a ograničenim iznad 40 N/mm.

\subsection{Razlikovanje tipa vrata}
\label{sec:tip-iz-sile}

Ishod ispitivanja razlikuje dva tipa vrata, kako prikazuje tablica \ref{tab:probe-osi}.

\begin{table}[H]
    \centering
    \caption{Očekivani ishod ispitivanja po osima prihvatnice}
    \label{tab:probe-osi}
    \begin{tabular}{@{}lll@{}}
        \toprule
        Os                         & Zakretna vrata & Klizna vrata \\
        \midrule
        prilazna, okomita na krilo & slobodna       & ograničena   \\
        os zatvaranja prstiju      & ograničena     & slobodna     \\
        uzduž kvake                & slobodna       & slobodna     \\
        \bottomrule
    \end{tabular}
\end{table}

Kod zakretnih vrata normala na krilo pri zatvorenom stanju poklapa se s tangentom luka između
šarke i kvake, pa je guranje duž prilazne osi ujedno i smjer otvaranja. Kod kliznih se vrata krilo
giba u vlastitoj ravnini, pa je slobodna os ona duž koje se prsti zatvaraju. Pravilo je time
jednostavno: koja je od tih dviju osi mekša, takva su vrata.

Os uzduž kvake ne koristi se jer je slobodna kod obaju tipova, budući da prihvatnica ondje klizi
duž same kvake. To je ujedno i načelno ograničenje metode: sila sama ne razlikuje gibanje vrata od
klizanja hvata, jer oboje daje nisku silu.

\subsection{Odnos prema razlikovanju iz percepcije}

Ispitivanje nije neovisan klasifikator. Da bi se uopće izvelo, prihvatnica mora držati kvaku, a za
hvat je potrebna pripremna poza koja ovisi o tipu kvake, okomitoj šipki naspram vodoravne poluge.
Tip se dakle mora pretpostaviti prije nego ispitivanje može početi.

Riječ je o dvostupanjskoj shemi. Prvi stupanj daje hipotezu iz percepcije, iz smjera vektora među
dvjema oznakama na kvaki (potpoglavlje \ref{sec:tip-kvake}). Drugi stupanj tu hipotezu provjerava
kroz interakciju, mjerenjem odziva sile. Vrijednost drugog stupnja nije u tome da zamijeni prvi,
nego da ga potvrdi mjerenjem koje ne ovisi o vidljivosti oznaka, osvjetljenju ni kutu gledanja, i
da usput izmjeri otpor, koji percepcija ne može dati.

U sustavu izvedenom u ovom radu razlikovanje tipa vrata provodi se isključivo iz
percepcije, bez ispitivanja silom. Razlog je taj što se prvi stupanj pokazao dovoljno pouzdanim.
U svim je pokretanjima tip vrata prepoznat točno, a ispitivanje silom tu je klasifikaciju potvrdilo
u svih dvadeset pokretanja, deset po tipu vrata (potpoglavlje \ref{sec:provjera}). Uvođenje drugog
stupnja u tijek izvođenja dodalo bi trajanje i kontaktne cikluse, a ne bi promijenilo nijednu
odluku.

Ispitivanje silom time u ovom radu ima ulogu neovisne provjere i mjernog postupka, a ne dijela
upravljačkog lanca. Njegova bi vrijednost došla do izražaja tek u okolini u kojoj oznake nisu
postavljene ili nisu pouzdano vidljive, gdje bi preuzelo ulogu prvog stupnja, a ne samo njegove
potvrde.

\section{Provjera procjena}
\label{sec:provjera}

Sve su procjene provjerene usporedbom sa stvarnim stanjem scene, koje simulator zna, ali koje u
upravljačku petlju nikada ne ulazi.

\subsection{Procjena osi šarke i kuta vrata}

Tablica \ref{tab:provjera-os} sažima odstupanja izmjerena kroz deset pokretanja otvaranja
zakretnih vrata, ukupno 72 koraka upravljanja.

\begin{table}[H]
    \centering
    \caption{Odstupanje procjene osi šarke i kuta vrata od stvarnih vrijednosti}
    \label{tab:provjera-os}
    \begin{tabular}{@{}lcccc@{}}
        \toprule
        Veličina                    & Sredina           & Medijan           & Najveća           & Std. devijacija   \\
        \midrule
        Pogreška položaja osi šarke & 5,8 mm            & 4,9 mm            & 14,3 mm           & 3,8 mm            \\
        Pogreška kuta vrata         & 0,28\textdegree{} & 0,23\textdegree{} & 1,46\textdegree{} & 0,28\textdegree{} \\
        \bottomrule
    \end{tabular}
\end{table}

Pogreška položaja osi raste s otvorenošću vrata, otprilike dvostruko između zatvorenog stanja i
65\textdegree{}. Pogreška kuta pritom nema takav trend, i razlog je geometrijski: pomak
procijenjene osi uglavnom je usmjeren prema kvaki, a kut se računa upravo iz tog smjera, pa
pomak duž njega u kut ne ulazi.

Procjena polumjera hvata pokazuje sustavno odstupanje od $-1{,}43$ mm, uz standardnu devijaciju od
svega 0,17 mm. Uska raspodjela oko konstantne vrijednosti pokazuje da je riječ o pristranosti, a ne
o šumu, tj. polumjer se dosljedno podcjenjuje.

\subsection{Razlikovanje tipa vrata i otpor}

Ispitivanje je provedeno u deset pokretanja po tipu vrata, pri čemu svako kreće od resetirane
scene i punog prilaska, pa raspršenje uključuje i promjenjivost samog hvata. Rezultati su u
tablici \ref{tab:provjera-probe}.

\begin{table}[H]
    \centering
    \caption{Ishod ispitivanja ograničenja, deset pokretanja po tipu vrata}
    \label{tab:provjera-probe}
    \begin{tabular}{@{}lcc@{}}
        \toprule
        Veličina                           & Zakretna  & Klizna    \\
        \midrule
        Točno razlikovanih                 & 10/10     & 10/10     \\
        Koeficijent, prilazna os, medijan  & 4,6 N/mm  & 99,8 N/mm \\
        Koeficijent, os prstiju, medijan   & 14,0 N/mm & 0,08 N/mm \\
        Omjer dvaju koeficijenata, medijan & 0,35      & 1081      \\
        \bottomrule
    \end{tabular}
\end{table}

Iz istog se ispitivanja dobiva i otporni moment zakretnih vrata, s medijanom od 260 Nm i rasponom
od 195 do 325 Nm. Ta se vrijednost poklapa s 280 do 308 Nm izračunatima iz sile izmjerene tijekom
samog otvaranja, dakle dvije neovisne metode daju isti otpor, što je najjača potvrda da
procjena mjeri svojstvo vrata, a ne artefakt postupka.

Treba ipak biti precizan oko toga što ta sila jest. Sila izmjerena u hvatu tijekom otvaranja nije
čisti otpor vrata, nego sadrži i silu kojom platforma gura kroz krutu vezu manipulatora. Zbog toga
se u ovom radu naziva silom u hvatu, a ne otporom vrata; iznos otpora dobiven ispitivanjem iz
mirovanja bliži je pravom svojstvu vrata.

\section{Uloga procjene u ovom radu}
\label{sec:procjena-uloga}

Dva se pristupa upravljanju bitno razlikuju po tome koliko se od gore opisanog stvarno koriste, i
to je jedna od glavnih razlika koja se razmatra u poglavlju \ref{ch:rezultati}.

Klasični pristup, opisan u poglavlju \ref{ch:klasicni}, geometriju ograničenja ne procjenjuje iz
gibanja, nego je uzima iz poze kvake dobivene percepcijom. Os šarke i smjer otvaranja poznati su
iz položaja i orijentacije hvatišta. Procjena sile i momenta ipak mu je nužna, za otkrivanje
uspješnog hvata i za ograničavanje kontaktne sile tijekom otvaranja. Ispitivanje ograničenja iz
potpoglavlja \ref{sec:constraint} u njegovu se izvođenju ne koristi, nego služi kao neovisna
provjera pretpostavke o tipu vrata i kao mjerenje otpora.

Pristup temeljen na učenju, opisan u poglavlju \ref{ch:rl}, geometriju ograničenja ne dobiva ni iz
percepcije ni iz eksplicitnog procjenitelja, nego ponašanje ograničenja mora zaključiti sam, iz
procijenjene sile i ostvarenog pomaka, i na temelju toga prilagoditi djelovanje. Time procjena sile
i momenta u tom pristupu nije pomoćni signal, nego jedini put kojim politika uopće doznaje o
mehaničkim svojstvima okoline.